%----------------------------------------------------------------------------------------
% NOTATIE EN FIGUREN VAN DE CURSUS STATISTIEK
% Staat in de preambule en niet in een hoofdstuk: ook de lijsten met oplossingen gebruiken dit.
%----------------------------------------------------------------------------------------
% Bewaarde grafiek, geplaatst door maak-pdf.py: bestand en pagina in figuren-cache/
\newcommand{\grafiekbestand}[2]{\includegraphics[page=#2,trim=3bp 3bp 3bp 3bp]{figuren-cache/#1}}
\usetikzlibrary{positioning}

% Combinaties, zoals in de cursus Telproblemen
\newcommand{\Comb}[2]{C_{#1}^{#2}}                         % combinaties
\newcommand{\N}{\mathbb{N}}

% Zoals stappen, maar met een bredere kolom voor de berekening: voor lange formules
\newenvironment{stappenbreed}{%
  \par\renewcommand{\arraystretch}{1.25}%
  \noindent\begin{tabular}{@{}>{$\displaystyle}p{9.8cm}<{$}@{\hspace{4mm}}>{\RaggedRight\small\itshape\arraybackslash}p{4.9cm}@{}}%
}{\end{tabular}\par}

% Zoals stappen, met eigen kolombreedtes: \begin{stappenvrij}{6.2cm}{3.8cm}. Voor naast een figuur.
\newenvironment{stappenvrij}[2]{%
  \par\renewcommand{\arraystretch}{1.25}%
  \noindent\begin{tabular}{@{}>{$\displaystyle}p{#1}<{$}@{\hspace{3mm}}>{\RaggedRight\small\itshape\arraybackslash}p{#2}@{}}%
}{\end{tabular}\par}

% Tekst of berekening links, figuur rechts: \naast[breedte rechts]{links}{rechts}
\newsavebox{\naastbox}
\newcommand{\naast}[3][]{%
  \par\noindent\sbox{\naastbox}{#3}%
  \begin{minipage}[t]{\dimexpr\linewidth-\wd\naastbox-3mm\relax}\vspace{0pt}#2\end{minipage}\hfill
  \begin{minipage}[t]{\wd\naastbox}\vspace{0pt}\usebox{\naastbox}\end{minipage}\par\addvspace{1mm}}
% Een rij figuren, gecentreerd
\newcommand{\figuren}[1]{\par\addvspace{1mm}{\centering #1\par}\addvspace{1.5mm}}
% Opgave met een grafiek rechts: \begin{metgrafiek}[breedte]{tikz} opgave en deelvragen \end{metgrafiek}
\newsavebox{\grafiekbox}
\newenvironment{metgrafiek}[2][]{%
  \sbox{\grafiekbox}{#2}%
  \begin{minipage}[t]{\dimexpr\linewidth-\wd\grafiekbox-3mm\relax}%
}{%
  \end{minipage}\hfill
  \begin{minipage}[t]{\wd\grafiekbox}\vspace{-2mm}\usebox{\grafiekbox}\end{minipage}\par}

% Gebeurtenissen en kansen
\newcommand{\compl}[1]{\overline{#1}}                    % complement van een gebeurtenis
\newcommand{\gegeven}{\,|\,}                              % voorwaardelijke kans: P(A \gegeven B)
\newcommand{\E}{\operatorname{E}}                         % verwachtingswaarde
\newcommand{\Var}{\operatorname{Var}}                     % variantie
\newcommand{\R}{\mathbb{R}}

% Boomdiagram, venndiagram en staafdiagram in de kleur van het hoofdstuk
\tikzset{
  tak/.style={accent,thick},
  blad/.style={font=\small,inner sep=1.5pt},
  kanslabel/.style={fill=white,inner sep=1pt,font=\small,text=black},
  verz/.style={draw=accent,thick,fill=accent,fill opacity=.10,text opacity=1},
}
% Staaf in een staafdiagram: \staaf{x}{hoogte}{breedte}
\newcommand{\staaf}[3]{\fill[accent!55] ({#1-#3/2},0) rectangle ({#1+#3/2},#2); \draw[accent] ({#1-#3/2},0) rectangle ({#1+#3/2},#2);}
% Compacte matrix in lopende tekst (zoals in Matrices): \km{1 & 2 \\ 3 & 4}
\newcommand{\km}[1]{{\setlength{\arraycolsep}{2.4pt}\renewcommand{\arraystretch}{.82}\left[\begin{array}{*{10}{c}}#1\end{array}\right]}}
% Dichtheid van de normale verdeling, voor plots: gauss(x, mu, sigma)
\tikzset{declare function={gauss(\x,\m,\s)=exp(-((\x-\m)^2)/(2*\s^2))/(\s*sqrt(2*pi));}}
\newcommand{\xbar}{\overline{x}}
